\documentclass[10pt]{beamer}
\usetheme{metropolis}
\usepackage{graphicx}
\usepackage{listings}
\usepackage{booktabs}
\usepackage{xcolor}
\usepackage{hyperref}

\definecolor{codebg}{HTML}{F5F5F5}
\definecolor{codegreen}{HTML}{2E7D32}
\definecolor{codegray}{HTML}{757575}
\definecolor{codeblue}{HTML}{1565C0}

\lstdefinestyle{rstyle}{
  language=R, backgroundcolor=\color{codebg}, basicstyle=\ttfamily\scriptsize,
  keywordstyle=\color{codeblue}\bfseries, stringstyle=\color{codegreen},
  commentstyle=\color{codegray}\itshape, breaklines=true, frame=single,
  rulecolor=\color{codegray!30}, showstringspaces=false, tabsize=2,
  morekeywords={library,ggplot,aes,geom_point,geom_line,geom_col,geom_histogram,
    geom_smooth,geom_boxplot,geom_density,
    facet_wrap,facet_grid,vars,scales,labeller,label_both,label_value,
    theme_minimal,theme,element_text,element_blank,strip.text,
    labs,coord_flip,scale_y_continuous,nrow,ncol}
}
\lstdefinestyle{pythonstyle}{
  language=Python, backgroundcolor=\color{codebg}, basicstyle=\ttfamily\scriptsize,
  keywordstyle=\color{codeblue}\bfseries, stringstyle=\color{codegreen},
  commentstyle=\color{codegray}\itshape, breaklines=true, frame=single,
  rulecolor=\color{codegray!30}, showstringspaces=false, tabsize=4,
  morekeywords={import,from,as,ggplot,aes,geom_point,geom_line,geom_histogram,
    geom_smooth,facet_wrap,facet_grid,theme_minimal,labs,
    FacetGrid,catplot,relplot,subplots,GridSpec}
}

\title{Faceting \& Small Multiples}
\subtitle{Workshop 2 --- Module 5}
\author{Prof.Asoc.\ Endri Raco}
\institute{Data Visualization for Data Scientists \\ UNYT Tirana}
\date{2025/26}

\begin{document}

\begin{frame}
  \titlepage
\end{frame}

\begin{frame}{Module Roadmap}
  \begin{enumerate}
    \item Tufte's small multiples revisited via the grammar
    \item \texttt{facet\_wrap()}: one variable, wrapped grid
    \item \texttt{facet\_grid()}: two variables, rows $\times$ columns
    \item Fixed vs free scales: the critical trade-off
    \item Facet vs colour: when to use which
    \item Faceting in Python: plotnine, seaborn FacetGrid, matplotlib
  \end{enumerate}
  \begin{exampleblock}{Learning Outcome}
    You will use faceting to produce Tufte-quality small multiples that
    separate groups into independent panels while preserving shared axes
    for fair comparison.
  \end{exampleblock}
\end{frame}

\section{Small Multiples via the Grammar}

\begin{frame}{From Spaghetti to Small Multiples}
  \begin{center}
    \includegraphics[width=0.92\textwidth]{figures_w02m05/facet_vs_color}
  \end{center}
  \begin{alertblock}{Pro-Tip}
    Colour works for $\leq$3 groups. For 4+ groups, faceting produces
    cleaner separation by exploiting Gestalt proximity (W01-M03)
    instead of hue discrimination.
  \end{alertblock}
\end{frame}

\begin{frame}{facet\_wrap: Concept}
  \begin{center}
    \includegraphics[width=0.85\textwidth]{figures_w02m05/facet_wrap_concept}
  \end{center}
\end{frame}

\section{facet\_wrap()}

\begin{frame}{facet\_wrap Demo: Six Regions}
  \begin{center}
    \includegraphics[width=0.92\textwidth]{figures_w02m05/facet_wrap_demo}
  \end{center}
\end{frame}

\begin{frame}[fragile]{facet\_wrap in R}
  \begin{columns}[T]
    \begin{column}{0.52\textwidth}
\begin{lstlisting}[style=rstyle]
ggplot(df, aes(x = month,
               y = revenue)) +
  geom_line(color = "#1565C0",
            linewidth = 1) +
  geom_area(fill = "#1565C0",
            alpha = 0.06) +

  # One panel per region
  facet_wrap(
    ~ region,
    ncol = 3,         # 3 columns
    scales = "fixed"  # shared y
  ) +

  scale_x_continuous(
    breaks = c(1, 6, 12),
    labels = c("Jan","Jun","Dec")) +
  theme_minimal(base_size = 9) +
  theme(strip.text =
    element_text(face = "bold")) +
  labs(title = "Revenue by Region",
       x = NULL, y = "Revenue ($K)")
\end{lstlisting}
    \end{column}
    \begin{column}{0.45\textwidth}
      \includegraphics[width=\textwidth]{figures_w02m05/r_facet}

      \vspace{4pt}
      {\small Key arguments:
      \texttt{ncol} / \texttt{nrow} control
      the grid layout.
      \texttt{scales = "fixed"} (default)
      ensures fair comparison.
      The strip text above each panel
      shows the faceting variable's level.}
    \end{column}
  \end{columns}
\end{frame}

\begin{frame}[fragile]{facet\_wrap in Python}
  \begin{columns}[T]
    \begin{column}{0.52\textwidth}
      \textbf{plotnine (ggplot2 clone)}
\begin{lstlisting}[style=pythonstyle]
from plotnine import *

(ggplot(df, aes(x="month",
                y="revenue"))
 + geom_line(color="#1565C0")
 + facet_wrap("~region", ncol=3)
 + theme_minimal())
\end{lstlisting}

      \vspace{6pt}
      \textbf{seaborn FacetGrid}
\begin{lstlisting}[style=pythonstyle]
import seaborn as sns

g = sns.FacetGrid(df,
    col="region", col_wrap=3,
    sharey=True, height=3)
g.map(plt.plot, "month","revenue",
      color="#1565C0")
g.set_titles("{col_name}")
\end{lstlisting}

      \vspace{6pt}
      \textbf{matplotlib (manual)}
\begin{lstlisting}[style=pythonstyle]
fig, axes = plt.subplots(
    2, 3, sharey=True)
for ax, region in zip(
    axes.flatten(), regions):
    sub = df.query(
        "region == @region")
    ax.plot(sub["month"],
            sub["revenue"])
    ax.set_title(region)
\end{lstlisting}
    \end{column}
    \begin{column}{0.45\textwidth}
      \includegraphics[width=\textwidth]{figures_w02m05/py_facet}

      \vspace{4pt}
      {\small Three Python approaches
      for faceting, ordered by
      automation level:
      \begin{enumerate}
        \item \texttt{plotnine}: identical
              to ggplot2 syntax
        \item \texttt{seaborn.FacetGrid}:
              map function to panels
        \item \texttt{plt.subplots}:
              full manual control
      \end{enumerate}}
    \end{column}
  \end{columns}
\end{frame}

\section{facet\_grid()}

\begin{frame}{facet\_grid: Rows $\times$ Columns}
  \begin{center}
    \includegraphics[width=0.92\textwidth]{figures_w02m05/facet_grid_demo}
  \end{center}
  \begin{exampleblock}{Data Insight}
    \texttt{facet\_grid(Type \textasciitilde{} Content\_Rating)} creates a
    matrix: rows = Type (Free, Paid), columns = Content Rating (Everyone,
    Teen, Mature). Every cell shows the same chart type for that
    combination. Use \texttt{facet\_grid} when you have two categorical
    variables; use \texttt{facet\_wrap} when you have one.
  \end{exampleblock}
\end{frame}

\begin{frame}[fragile]{facet\_grid in Code}
  \begin{columns}[T]
    \begin{column}{0.48\textwidth}
      \textbf{R}
\begin{lstlisting}[style=rstyle]
ggplot(apps, aes(x = Rating)) +
  geom_histogram(bins = 20,
    fill = "#1565C0",
    color = "white") +

  # rows = Type, cols = Content Rating
  facet_grid(
    Type ~ `Content Rating`,
    scales = "fixed") +

  theme_minimal(base_size = 9) +
  labs(title = "Rating by Type
       x Content Rating")
\end{lstlisting}
    \end{column}
    \begin{column}{0.48\textwidth}
      \textbf{Python (plotnine)}
\begin{lstlisting}[style=pythonstyle]
(ggplot(apps, aes(x="Rating"))
 + geom_histogram(bins=20,
     fill="#1565C0", color="white")
 + facet_grid("Type~Content_Rating")
 + theme_minimal())
\end{lstlisting}

      \vspace{4pt}
      \textbf{seaborn}
\begin{lstlisting}[style=pythonstyle]
g = sns.FacetGrid(apps,
    row="Type",
    col="Content_Rating",
    sharey=True, sharex=True)
g.map(plt.hist, "Rating",
      bins=20, color="#1565C0")
\end{lstlisting}
    \end{column}
  \end{columns}
\end{frame}

\section{Fixed vs Free Scales}

\begin{frame}{The Critical Trade-Off}
  \begin{center}
    \includegraphics[width=0.92\textwidth]{figures_w02m05/scales_fixed_vs_free}
  \end{center}
\end{frame}

\begin{frame}{Scale Options}
  \centering
  \begin{tabular}{lp{5.5cm}}
    \toprule
    \textbf{Setting} & \textbf{Effect} \\
    \midrule
    \texttt{scales = "fixed"} & All panels share x and y (default). Best for fair comparison across groups. \\
    \texttt{scales = "free\_x"} & Each panel has its own x range. Y stays shared. \\
    \texttt{scales = "free\_y"} & Each panel has its own y range. X stays shared. \\
    \texttt{scales = "free"} & Both axes independent. Best for showing within-group detail. \\
    \bottomrule
  \end{tabular}

  \vspace{6pt}
  \begin{alertblock}{Pro-Tip}
    \textbf{Default to fixed scales} for comparison tasks. Use free scales
    only when groups have vastly different ranges (e.g., countries with
    populations from 1K to 1B). When using free scales, warn the reader
    with an annotation: ``Note: axes differ across panels.''
  \end{alertblock}
\end{frame}

\section{Facet + Geom Combinations}

\begin{frame}{Facet + Smooth: Per-Panel Trends}
  \begin{center}
    \includegraphics[width=0.95\textwidth]{figures_w02m05/facet_smooth}
  \end{center}
  \begin{exampleblock}{Data Insight}
    When \texttt{geom\_smooth()} is combined with \texttt{facet\_wrap()},
    the smoother is fitted independently within each panel. This
    reveals group-specific trends that would be averaged out in a
    single-panel, colour-coded plot.
  \end{exampleblock}
\end{frame}

\begin{frame}{Custom Strip Labels}
  \begin{center}
    \includegraphics[width=0.92\textwidth]{figures_w02m05/strip_labels}
  \end{center}
  \begin{alertblock}{Pro-Tip}
    Use \texttt{labeller = label\_both} to show ``Variable: Level''
    or a custom labeller function to add sample sizes. In matplotlib,
    use \texttt{ax.set\_title(f"\{name\} (n=\{count\})")}.
  \end{alertblock}
\end{frame}

\section{Facet vs Patchwork}

\begin{frame}{When to Facet vs When to Compose}
  \begin{center}
    \includegraphics[width=0.88\textwidth]{figures_w02m05/patchwork_vs_facet}
  \end{center}
\end{frame}

\begin{frame}[fragile]{Decision Rule in Code}
\begin{lstlisting}[style=rstyle]
# FACET: Same geom, same aes, different group
ggplot(df, aes(x = month, y = revenue)) +
  geom_line() +
  facet_wrap(~region, ncol = 3)

# PATCHWORK: Different geom or different aes per panel
library(patchwork)
p_bar <- ggplot(df, aes(x = cat, y = n)) + geom_col()
p_hist <- ggplot(df, aes(x = rating)) + geom_histogram()
p_scatter <- ggplot(df, aes(x = reviews, y = rating)) + geom_point()
p_box <- ggplot(df, aes(x = type, y = rating)) + geom_boxplot()

# Compose: | = side by side, / = stack
(p_bar | p_hist) / (p_scatter | p_box) +
  plot_annotation(tag_levels = "a")
\end{lstlisting}
\end{frame}

\section{Synthesis}

\begin{frame}{Faceting Decision Checklist}
  \begin{enumerate}
    \item \textbf{How many groups?} 1--3: colour. 4+: facet.
    \item \textbf{How many faceting variables?} 1: \texttt{facet\_wrap}.
          2: \texttt{facet\_grid}.
    \item \textbf{Comparable ranges?} Yes: \texttt{scales = "fixed"}.
          No: \texttt{scales = "free"} + warning label.
    \item \textbf{Same geom across panels?} Yes: facet.
          No: patchwork / GridSpec.
    \item \textbf{Need independent smoothers?} Yes: facet +
          \texttt{geom\_smooth()} fits per panel automatically.
    \item \textbf{Need context in strip labels?} Use
          \texttt{labeller = label\_both} or custom function.
  \end{enumerate}
\end{frame}

\begin{frame}{Key Takeaways}
  \begin{enumerate}
    \item \textbf{facet\_wrap(\textasciitilde{}var)} splits by one variable into a
          wrapped grid; \textbf{facet\_grid(row\textasciitilde{}col)} splits by two
          variables into a matrix.
    \item \textbf{Fixed scales} (default) enable fair comparison;
          \textbf{free scales} reveal within-group detail at the cost of
          cross-panel comparability.
    \item Faceting produces Tufte's \textbf{small multiples} automatically:
          same chart, same axes, different data subset per panel.
    \item \textbf{Colour} for $\leq$3 groups; \textbf{facet} for 4+.
    \item \textbf{Facet} when all panels use the same geom;
          \textbf{patchwork/GridSpec} when panels use different geoms.
    \item In Python: \texttt{plotnine.facet\_wrap}, \texttt{sns.FacetGrid},
          or \texttt{plt.subplots} with a loop.
  \end{enumerate}
\end{frame}

\begin{frame}{Essential Readings}
  \begin{itemize}
    \item Wickham, H. (2016). \emph{ggplot2}, Chapter 7 (Facets).
    \item Tufte, E.~R. (1983). \emph{The Visual Display of Quantitative
          Information}, pp.\ 170--174 (Small Multiples).
    \item Wilke, C.~O. (2019). \emph{Fundamentals of Data Visualization},
          Chapter 21 (Multi-Panel Figures).
    \item Wickham, H. \& Grolemund, G. (2023). \emph{R for Data Science},
          2nd ed., Chapter 9.5 (Facets).
  \end{itemize}
  \vspace{6pt}
  \textbf{Next Module}: Themes \& Publication-Quality Styling (W02-M06)
\end{frame}

\end{document}
